% Supplement C: Spatial Structure Analysis
% Detailed analysis of minority clustering and Moran's I

Our finding of method equivalence depends critically on the spatial structure of minority populations in our test states. We analyze this structure using spatial autocorrelation metrics and discuss when our results generalize.

\subsection{Spatial Autocorrelation: Moran's I}

We quantify minority clustering using Moran's I~\cite{moran1950notes}, a standard measure of spatial autocorrelation. For a spatial variable $x$ (minority VAP percentage) and spatial weights matrix $W$ (tract adjacency):

\[
I = \frac{n}{\sum_i \sum_j w_{ij}} \cdot \frac{\sum_i \sum_j w_{ij}(x_i - \bar{x})(x_j - \bar{x})}{\sum_i (x_i - \bar{x})^2}
\]

where $n$ is the number of tracts, $w_{ij} = 1$ if tracts $i$ and $j$ are adjacent, and $\bar{x}$ is the mean minority percentage.

\textbf{Interpretation}:
\begin{itemize}
    \item $I > 0$: Positive spatial autocorrelation (clustering) --- minority-dense tracts are spatially adjacent
    \item $I \approx 0$: No spatial pattern (random distribution)
    \item $I < 0$: Negative spatial autocorrelation (dispersion) --- minority-dense tracts are spatially separated
\end{itemize}

\subsection{Spatial Clustering in Test States}

Table~\ref{tab:spatial-autocorrelation-supp} shows Moran's I for our five test states:

\begin{table}[h]
\centering
\begin{tabular}{lcccc}
\toprule
\textbf{State} & \textbf{Moran's I} & \textbf{z-score} & \textbf{p-value} & \textbf{Interpretation} \\
\midrule
Alabama & 0.712 & 45.3 & < 0.001 & Strong clustering \\
Georgia & 0.683 & 52.1 & < 0.001 & Strong clustering \\
Louisiana & 0.698 & 38.7 & < 0.001 & Strong clustering \\
Mississippi & 0.745 & 29.4 & < 0.001 & Very strong clustering \\
South Carolina & 0.679 & 41.2 & < 0.001 & Strong clustering \\
\midrule
\textbf{Average} & \textbf{0.703} & \textbf{41.3} & --- & \textbf{Strong clustering} \\
\bottomrule
\end{tabular}
\caption{Spatial autocorrelation (Moran's I) for minority VAP distribution in test states. All states exhibit statistically significant positive spatial autocorrelation (p < 0.001), indicating strong clustering of minority-dense census tracts. High clustering enables method equivalence by creating well-defined spatial clusters that all algorithms identify identically.}
\label{tab:spatial-autocorrelation-supp}
\end{table}

\textbf{Key findings}:
\begin{enumerate}
    \item All five states show strong positive spatial autocorrelation ($I > 0.65$)
    \item Mississippi exhibits the highest clustering ($I = 0.745$), consistent with its compact minority geography (Delta region)
    \item Statistical significance is overwhelming ($z > 29$ for all states), ruling out random spatial patterns
\end{enumerate}

\subsection{Relationship to Method Equivalence}

Spatial clustering directly enables method equivalence through the following mechanism:

\begin{enumerate}
    \item \textbf{Cluster formation}: High Moran's I means minority-dense tracts form contiguous spatial clusters

    \item \textbf{Edge-weighting amplification}: When $\alpha \gg 1$, edges within clusters receive weight $\alpha$, while edges crossing cluster boundaries receive weight 1

    \item \textbf{Optimization convergence}: All partitioning algorithms solve the same problem: ``partition the state without cutting high-weight edges.'' The spatial clustering makes this problem have a unique (or near-unique) solution.

    \item \textbf{Method irrelevance}: Since the cluster structure is obvious (high Moran's I), all algorithms identify it identically, making methodological choice irrelevant
\end{enumerate}

\textbf{Formal connection}: Recall from the main paper that method equivalence requires $\alpha \geq \alpha_{\text{crit}} \approx k / m_{\text{state}}$. However, this bound assumes worst-case spatial distribution. With high spatial autocorrelation, the effective $\alpha_{\text{crit}}$ is lower:

\[
\alpha_{\text{crit}}^{\text{effective}} = \frac{k}{m_{\text{state}} \cdot I}
\]

where $I$ is Moran's I. This explains why our empirical results show convergence at $\alpha = 5$ despite theoretical predictions of $\alpha_{\text{crit}} \in [10, 38]$: spatial clustering reduces the threshold.

\subsection{Generalization Boundaries}

Our results generalize to states with high minority clustering but may not hold for dispersed minority populations:

\begin{table}[h]
\centering
\begin{tabular}{lccl}
\toprule
\textbf{Scenario} & \textbf{Moran's I} & \textbf{Predicted Behavior} & \textbf{Method Dependence?} \\
\midrule
High clustering & $I > 0.6$ & Equivalence at $\alpha \approx 5$ & No \\
Moderate clustering & $0.3 < I < 0.6$ & Equivalence at $\alpha \approx 20$ & Low \\
Low clustering & $0.1 < I < 0.3$ & Equivalence at $\alpha \approx 50$ & Moderate \\
Random/dispersed & $I < 0.1$ & Equivalence at $\alpha \gg 100$ & High \\
\bottomrule
\end{tabular}
\caption{Predicted method equivalence thresholds as a function of spatial clustering (Moran's I). States with dispersed minority populations require higher $\alpha$ to achieve method equivalence.}
\label{tab:generalization-supp}
\end{table}

\textbf{Case studies}:
\begin{itemize}
    \item \textbf{High clustering states} (Black Belt South, Hispanic Southwest): Our results apply directly. Examples: Alabama, Mississippi, Louisiana (this study), Texas, New Mexico.

    \item \textbf{Moderate clustering states} (Urban-concentrated minorities): Higher $\alpha$ needed but equivalence still achievable. Examples: Illinois (Chicago), California (Los Angeles).

    \item \textbf{Low clustering states} (Dispersed rural minorities): Very high $\alpha$ needed; method choice may remain relevant. Examples: Montana, Wyoming, Vermont (small minority populations distributed across rural areas).
\end{itemize}

\subsection{Implications for Redistricting Practice}

Spatial structure analysis has practical implications:

\begin{enumerate}
    \item \textbf{Pre-assessment}: Before running partitioning algorithms, compute Moran's I for minority distribution. High I ($> 0.6$) suggests method equivalence will hold; low I suggests testing multiple methods.

    \item \textbf{Alpha calibration}: Use spatial autocorrelation to calibrate edge weight factor:
    \[
    \alpha_{\text{recommended}} = \frac{2k}{m_{\text{state}} \cdot \max(I, 0.5)}
    \]
    This ensures robustness across varying spatial structures.

    \item \textbf{Litigation strategy}: In VRA cases, plaintiffs can use Moran's I to argue that method equivalence should hold, reducing discretion in algorithm choice. Conversely, defendants facing dispersed minorities ($I < 0.3$) can argue that method choice legitimately affects outcomes.

    \item \textbf{Sensitivity testing}: States with moderate I ($0.3 < I < 0.6$) should test robustness by running multiple algorithms and verifying convergence.
\end{enumerate}

\subsection{Open Questions}

Several spatial structure questions remain for future research:

\begin{itemize}
    \item \textbf{Temporal dynamics}: How does Moran's I evolve across census decades as demographics shift? Does increasing suburban diversity reduce clustering and weaken method equivalence?

    \item \textbf{Scale dependence}: Does clustering at tract level ($n \approx 1000$) translate to clustering at block group level ($n \approx 5000$) or precinct level ($n \approx 2000$)?

    \item \textbf{Multi-group clustering}: We analyze Black+Hispanic minorities aggregated. Do individual racial groups show different clustering patterns, and does this affect method equivalence for single-group VRA compliance?

    \item \textbf{Alternative metrics}: Moran's I captures first-order autocorrelation. Do higher-order statistics (Geary's C, LISA) reveal additional structure affecting method equivalence?
\end{itemize}

\subsection{Summary}

Our test states exhibit strong minority clustering (average Moran's I = 0.703), creating well-defined spatial structure that all partitioning algorithms identify identically when using edge-weighting with $\alpha \geq 5$. This spatial structure is the fundamental driver of method equivalence.

The finding generalizes to states with similar clustering (I > 0.6) but may require higher $\alpha$ for dispersed minority populations (I < 0.3). Practitioners should assess spatial structure before algorithm selection to determine appropriate edge weight calibration.
